\section{Introduction}
\label{sec:intro}

Applied problems often present an effect that is not constant across the
sample. The return to schooling changes with age and with experience, the
effect of temperature on electricity demand changes with the hour of the day
and with humidity, and the effect of a treatment changes with the age and the
body mass index of the patient. The varying coefficient model of
\citet{Hastie-Tibshirani-1993} and \citet*{Cai-Fan-Yao-2000} accommodates this
by letting the coefficient of each covariate be a function of a modulating
variable; see \citet{Fan-Zhang-2008} for a review. When more than one variable
modulates the effect, a fully nonparametric coefficient runs into the curse of
dimensionality, and the standard remedy is to make each coefficient additive
in the modulators, which is the additive coefficient model of
\citet{Xue-Yang-2006}, studied further by \citet{Xue-Liang-2010} and
\citet{Liu-Yang-2010}.

The estimators available for this model are built from polynomial splines and
from kernels, and are governed by a global smoothness index, so that a
component with a peak, a jump or a change of regime is either oversmoothed or
paid for everywhere in the domain. Wavelet bases combined with an $\ell_1$
penalty \citep{tibshirani1996regression} are the standard answer to this difficulty, because a function whose
smoothness is inhomogeneous is still economically represented by a few large
coefficients \citep{Donoho-Johnstone-1994, Donoho-Johnstone-1998,
Antoniadis-Fan-2001}. They have been used for varying coefficients without an
additive structure and without penalization \citep*{Zhou-You-2004,
Montoril-Morettin-Chiann-2018}, and for additive models without a multiplying
covariate \citep{Sardy-Tseng-2004, Zhang-Wong-2003, Haris-Simon-Shojaie-2018,
Amato-Antoniadis-DeFeis-Gijbels-2022, Sardy-Ma-2024}; penalized estimation in
the additive coefficient model itself has been carried out with grouped
penalties on spline bases \citep{Wei-Huang-Li-2011,
Antoniadis-Gijbels-LambertLacroix-2014, Ma-Carroll-Liang-Xu-2015}, again under
global smoothness.

\citet{Klopp-Pensky-2015} study the varying coefficient model in which a
single index modulates every coefficient. They expand each coefficient
function in an orthonormal basis, estimate the resulting array by a block
lasso, and obtain a nonasymptotic oracle inequality, a Besov rate that adapts
to inhomogeneous smoothness, and a matching minimax lower bound, under the
assumption that the covariates are independent of the index. The present paper
carries that  program to the situation an applied problem usually presents,
in which several variables modulate a coefficient at once. We let each
coefficient be additive in the modulators, which keeps the estimator free of
the curse of dimensionality in their number, we allow the covariates to depend
on the modulators, and we  keep the block lasso but leave the levels of the
coefficients unpenalized and group the wavelet coefficients in chunks of order
$\log n$ within a resolution level, with weights whose ratio stays bounded, so
that the standardized weights of group lasso software are covered.
Additivity produces cross terms between distinct modulators,
so the population Gram matrix no longer factorizes as a Kronecker product and
has to be bounded directly; dependence between the covariates and the
modulators replaces a marginal second moment by a conditional one; and the
unpenalized levels leave a term in the risk that full block penalization does not
incur. Our main result, Theorem~\ref{thm:main}, states that the estimator
attains the rate of  
Klopp and Pensky, with no logarithmic factor when the Besov integrability
index is at least two, 
% E5g (D57): "where it is minimax optimal", pela cota inferior do Theorem S6.1.
where it is minimax optimal. It holds under weaker conditions than
theirs: the number of covariates is fixed, a regime their high dimensional
condition excludes, and the effective smoothness of the components need only
exceed $s/(2s+1)$, which is below one half, where theirs must exceed one half.
Components with effective smoothness one half, such as the inhomogeneous
functions of Section~\ref{sec:simulation}, are covered here and not there. 
% E5g (D57): a frase de D50 sobre a cota inferior, trocada (Theorem S6.1 do supp).
% D58 (pergunta 52(d)): "it" trocado por "the rate".
The lower bound that makes  the rate optimal is proved here for any number of
modulators and for covariates that depend on them; Klopp and Pensky have it
with a single modulator independent of the covariates.
% E5h (k = 5): o "Section~5" acima passou a \ref{sec:simulation} (D50); o texto
% impresso nao muda. A Secao 5 diz que as funcoes nao homogeneas tem s' = 1/2.

We call the resulting procedure WAFC, for wavelet additive functional
coefficients. Its contributions are of   three kinds. The first is the theory of
the additive product design: a design condition for the Gram matrix of the
products between a linear covariate and a wavelet of a modulator, which
carries a cross term between distinct modulators and does not assume
independence (Proposition~\ref{prop:gram}); an oracle inequality that needs no
cone condition, because the unpenalized levels are profiled out, and
that holds for any weights of bounded ratio, including those of the software
that computes the estimator
(Theorem~\ref{thm:oracle}); and rates 
that follow from the Besov assumption alone, through a bound on the
ideal risk over chunks
(Theorem~\ref{thm:main}), together with a minimax lower bound, for any number
of modulators, that makes these rates optimal when $\pi \ge 2$ (Theorem~S6.1 in
the Supplementary Material). 
The second is the recovery of the structure of the model, which
modulators act on which coefficients, by a threshold on the norms of the
fitted components (Corollary~\ref{cor:threshold}): removing the null blocks
needs no separation condition, and the thresholded estimator keeps the rate of
the fit (Corollary~\ref{cor:thrisk}), so that one estimator serves both for
structure and for prediction. The third is computational: the
estimator is one convex program  solved by
standard group lasso software,  the resolution
level and the penalty are chosen jointly by cross-validation, and the
threshold by the one standard error rule on the same folds.


The rest of the paper is organized as follows. Section~\ref{sec:model}
introduces the model, its identification, the wavelet approximation space and the estimator.
Section~\ref{sec:theory} states the assumptions and the theory, opening with
the main rate and then giving the three results it rests on, then
the rates, including one with no design condition, and
closes with the recovery of the block structure by thresholding and
the risk of the thresholded estimator.
Section~\ref{sec:computation} describes the computation, the scale of the
penalty and the tuning rules, the threshold among them.
% E5h (k = 5): as Secoes 5 e 6 no roteiro; a Secao 7 (Discussion) e a frase de
% fechamento continuam com a E5b.
Section~\ref{sec:simulation} compares the method with penalized
splines and with other competitors in a simulation study, and
Section~\ref{sec:application} applies it to hourly air pollution data from
London. Section~\ref{sec:discussion} discusses where the method
does not win and what is left open. Proofs are collected in the online Supplementary
Material, which also gives the theory of the coordinatewise lasso,
offered by the software as an option, the details and further
results of the simulation study, and a second application, to data from
Beijing.

Throughout, $\norm{a}_\pi$ is the $\ell_\pi$ norm of a vector,
$\normn{h}^2 = n^{-1}\sum_{i\le n} h_i^2$ is the empirical norm,
$\lmin(\cdot)$ and $\lmaxx(\cdot)$ are the extreme eigenvalues of a symmetric
matrix, $\opnorm{\cdot}$ is the operator norm, and $a_n \asymp b_n$ means that
$a_n/b_n$ is bounded away from zero and infinity. For a support
$S$, $\delta_S$ is the vector that agrees with $\delta$ on $S$ and vanishes
elsewhere.

%%%%%%%%%%%%%%%%%%%%%%%%%%%%%%%%%%%%%%%%%%%%%%%%%%%%%%%%%%%%%%%%%%%%%%%%%%%%%%
